% Lane WA — design decisions + deviations. Uses longtable/booktabs (preamble).
\section{Design Decisions}\label{sec:decisions}

TensixFlow was planned before any code was written, and the plan kept an explicit
decision register (drafts D1--D32 in the plan workspace) that was revised by three
rounds of adversarial review. Table~\ref{tab:decisions-register} records the
decisions that shaped the delivered system; Section~\ref{sec:decisions-deviations}
records the points where execution deliberately and documentedly deviated from the
plan text. The numerical and architectural background referenced in the register is
developed in Section~\ref{sec:theory}.

\subsection{Decision register}

{\small
\setlength{\tabcolsep}{4pt}
\begin{longtable}{@{}p{1.1cm}p{3.55cm}p{2.9cm}p{4.75cm}p{1.5cm}@{}}
\caption{Decision register. IDs are the plan's draft-decision identifiers.
  ``Adopted'' means the decision was made during planning and the delivered system
  implements it; revisions note the review round that forced them.}
\label{tab:decisions-register}\\
\toprule
\textbf{ID} & \textbf{Decision} & \textbf{Alternatives considered}
  & \textbf{Rationale} & \textbf{Status} \\
\midrule
\endfirsthead
\multicolumn{5}{@{}l}{\footnotesize\textit{Table \ref{tab:decisions-register}
  (decision register) continued}}\\
\toprule
\textbf{ID} & \textbf{Decision} & \textbf{Alternatives considered}
  & \textbf{Rationale} & \textbf{Status} \\
\midrule
\endhead
\midrule
\multicolumn{5}{r@{}}{\footnotesize\textit{continued on next page}}\\
\endfoot
\bottomrule
\endlastfoot

D1
& Self-contained, out-of-tree compiler: own \texttt{tensixflow} MLIR dialect and a
  full lowering pipeline living in this repository.
& Building in-tree inside Tenstorrent's tt-forge / tt-mlir repositories.
& The only way to make hardware-free progress: the standalone stack builds and
  tests on an ordinary CPU machine with the nix LLVM 20.1.8 toolchain, needs no
  Tenstorrent checkout, and keeps every artifact inspectable. Compatibility with
  the Tenstorrent stack is handled later as a TTIR emitter plus a prose mapping
  doc, not by living inside their tree.
& Adopted \\

D18
& Toolchain spike first: build the upstream \texttt{llvm-project}
  \texttt{mlir/examples/standalone} example under the new nix flake + devenv
  shell before writing any dialect code.
& Starting dialect work directly against the new environment.
& The devenv + MLIR wiring (\texttt{MLIR\_DIR}/\texttt{LLVM\_DIR}/tblgen paths)
  was explicitly UNVERIFIED at plan time. A throwaway spike proves the wiring
  empirically; had it failed, the environment approach would have been revisited
  before any dialect investment.
& Adopted; spike passed \\

D22
& Pipeline-registration probe: prove \texttt{registerAllPasses()} plus multi-pass
  composition with a 2-pass lit test before relying on the full seven-pass
  pipeline.
& Building the full pipeline first and debugging registration through it.
& The synthesized pipeline-registration syntax was UNVERIFIED at plan time; the
  2-pass probe isolates registration risk from pass-logic risk.
& Adopted; probe passed \\

D14 (execution)
& Numerical validation at the schedule-IR level: the functional simulator executes
  the schedule IR produced by \texttt{-{}-tensixflow-lower-to-ttkernel} over
  modeled circular buffers.
& Executing the generated kernels on a device (impossible: no device owned);
  interpreting the emitted C++ (out of scope).
& No Tenstorrent device exists in this project, so the generated
  reader/compute/writer/host sources are validated structurally only (golden
  FileCheck against current-API tokens + compile-only check) and are never
  executed. Keeping the numeric claim one level above the codegen artifact is what
  keeps the honesty disclaimer true.
& Adopted \\

D14 (rounding)
& Deliberate precision gap: simulator = bf16 spatial model (inputs bf16, fp32
  accumulate within a tile, per-tile bf16 rounding, bf16 circular-buffer
  intermediates); golden = fp32-exact reference. See
  Figure~\ref{fig:theory-precision-flow} and Eq.~\eqref{eq:theory-pcc}.
& Original contract --- golden also rounds per-op to bf16 (rejected in review
  round 1, P0); a hard-coded $0.9999$ upper guard (rejected in review round 2,
  R1).
& Revised per adversarial review so the PCC measures real divergence: the original
  contract collapsed PCC $\to 1.0$ and made the 0.9875 bar vacuous; the hard-coded
  upper guard was miscalibrated and would false-fail a correct implementation
  (good-case PCC exceeds it for single ops). The delivered contract makes
  PCC(simulator\textsubscript{bf16 tiled}, golden\textsubscript{fp32 exact}) the
  measured cost of fusion + tiled execution, with non-vacuity guarded by a
  structural dtype assertion, a bit-difference check, and an empirically calibrated
  trip-wire.
& Adopted after two review revisions \\

D4
& Exactly one modeled device: Wormhole n150 class, 8x8 grid, 32x32 tiles,
  bf16/fp32, HiFi4, with a modeled 1\,MiB L1 per core.
& A multi-device target matrix; using the literal silicon budget.
& A single target avoids premature generalization, and the modeled 1\,MiB is a
  documented, conservative, tunable constant: the literal n150 L1 per core is
  larger ($\approx$1.5\,MB), so the planner stays safely inside the real budget.
  Changing the one constant (and re-checking planner + model) is the supported way
  to target another budget.
& Adopted \\

D3
& Codegen targets the \emph{current} tt-metal API at \texttt{@df36cd6}: single
  \texttt{CreateKernel} with a \texttt{variant} of
  \texttt{DataMovementConfig}/\allowbreak\texttt{ComputeConfig}, \texttt{tt::CBIndex}
  circular buffers, \texttt{distributed::Enqueue*} mesh API.
& Emitting the legacy \texttt{EnqueueProgram} / per-role \texttt{CreateKernel}
  surface.
& \texttt{EnqueueProgram} no longer exists in the current \texttt{host\_api.hpp};
  it is kept only as a mapping note in \texttt{docs/tenstorrent\_mapping.md}. The
  golden structural tests assert current-API tokens only, and a legacy-API fixture
  is deliberately kept failing so the stub set cannot silently drift back to the
  legacy surface.
& Adopted \\

D15
& A single JSON schema as the plan/perf/viz contract, owned before the consumers
  and imported by all of them.
& Separate ad-hoc schemas per consumer (scheduler, perf model, report tool, HTML
  visualizer).
& One schema keeps the scheduler's plan, the perf model's figures, and the
  visualizer's sections from drifting apart. It is deliberately \emph{separate}
  from the MLIR-level schedule IR (D14), which is the codegen$\leftrightarrow$simulator
  contract; the two contracts serve different consumers.
& Adopted \\

D16
& Signature-only tt-metal stub headers as the substrate for the compile-only
  kernel check.
& Building real tt-metal (its runtime requires a device); skipping the compile
  check entirely.
& With no device, the stubs --- pinned to \texttt{@df36cd6} and covering both the
  host-API surface and the kernel-side intrinsics --- are the strongest available
  compile-time statement about the emitted code. The check asserts the presence of
  expected emitted symbols so it cannot pass vacuously, and a negative legacy-API
  fixture proves the stub set rejects the removed surface.
& Adopted \\

D17
& Shared L1-budget model: $\text{CB+shard budget} = L1_{\text{total}} -
  \text{reserve}$, $\text{CB\_bytes} = n_{\text{CB}} \times \text{depth} \times
  32 \times 32 \times 2$, defined once in \texttt{L1Model.h} and imported verbatim
  by the circular-buffer allocator and the perf model.
& Each component carrying its own copy of the constants.
& Single source of truth: the occupancy the perf model reports is exactly the
  budget the allocator actually planned against (Figure~\ref{fig:theory-core-memory}).
  Duplicating the constants would let the two silently disagree.
& Adopted \\

\end{longtable}
}

\subsection{Deviation register}\label{sec:decisions-deviations}

Four points where execution deviated from the literal plan text. Each deviation was
documented at the time it was taken; none of them loosened a verification gate.

{\small
\setlength{\tabcolsep}{4pt}
\begin{longtable}{@{}p{0.9cm}p{3.5cm}p{8.4cm}p{2.0cm}@{}}
\caption{Deviation register. ``T<number>'' refers to the plan's todo number whose
  acceptance prose was affected.}
\label{tab:decisions-deviation-register}\\
\toprule
\textbf{ID} & \textbf{Plan said} & \textbf{Deviation taken and why}
  & \textbf{Status} \\
\midrule
\endfirsthead
\multicolumn{4}{@{}l}{\footnotesize\textit{Table
  \ref{tab:decisions-deviation-register} (deviation register) continued}}\\
\toprule
\textbf{ID} & \textbf{Plan said} & \textbf{Deviation taken and why}
  & \textbf{Status} \\
\midrule
\endhead
\bottomrule
\endlastfoot

T37
& Acceptance prose invoked \texttt{-{}-emit-metalium \ldots -o <dir>} to write the
  generated project.
& A bare \texttt{-o <dir>} is impossible for an MLIR pass option:
  \texttt{MlirOptMain} owns the global \texttt{-o} output-file flag, and
  \texttt{PassOptions} options are parsed only out of the pass flag's value string.
  Implemented and documented as
  \texttt{-{}-tensixflow-emit-metalium=output-dir=<dir>} (recorded in the pass's
  TableGen description and the lit test header); acceptance was met in the
  deviated form.
& Documented deviation \\

T50
& Optional tt-mlir build + \texttt{ttmlir-opt} study, off the critical path.
& The heavy build (tt-mlir compiles a private 30--60\,min LLVM toolchain) was
  skipped on this box; the plan-sanctioned disjunctive OR branch was taken instead:
  \texttt{docs/ttmlir\_study.md} documents the skip reason and carries verbatim
  source-study snippets from tt-mlir's own test files at \texttt{@9f06802f}.
& Passed via OR branch \\

T55
& On-device branch: compile the generated kernels against real tt-metal, run them,
  and validate PCC $\geq 0.9875$ against the golden.
& Skipped: there is no \texttt{/dev/tenstorrent} on this machine, so the branch is
  gated out by design. Acceptance is covered by the disjunctive hardware harness,
  which passes via its explicit if/else no-device branch (bring-up dry-run +
  compile-only kernel check) --- the on-device branch was never silently replaced,
  only deferred.
& Skipped (gated); not failed \\

T60
& Submit the prepared upstream contribution as a CI-green PR (or patch + issue).
& Submission skipped per the explicit owner instruction ``DO NOT SUBMIT ANY PRS''.
  The contribution was implemented and verified locally (upstream pre-commit suite
  green, checks run against the current upstream main) and is committed to the
  repository's \texttt{upstream/} tree as a documentation-only artifact; no fork,
  push, PR, or issue was created.
& Documented only, per owner instruction \\

\end{longtable}
}

Two conventions close the register. Statuses are honest about what was \emph{not}
done: a gated skip is recorded as a skip, and a fallback branch is recorded as the
fallback, never as the primary outcome. And the honesty boundary from
Section~\ref{sec:theory} is preserved verbatim in every row: nothing in the
register upgrades ``structurally validated, never executed'' to ``executed'' or
``numerically verified''.
